% Generated draft: requires mathematical and source review.
\documentclass[11pt]{article}
\usepackage[margin=25mm]{geometry}
\usepackage{fontspec}
\setmainfont{Latin Modern Roman}
\usepackage{amsmath,amssymb,mathtools,unicode-math}
\setmathfont{Latin Modern Math}
\usepackage{enumitem,needspace,xurl,hyperref}
\setlength{\parindent}{0pt}
\setlength{\parskip}{5pt}
\setlength{\emergencystretch}{3em}
\newcommand{\N}{\mathbb{N}}
\newcommand{\Z}{\mathbb{Z}}
\newcommand{\Q}{\mathbb{Q}}
\newcommand{\R}{\mathbb{R}}
\newcommand{\C}{\mathbb{C}}
\newcommand{\F}{\mathbb{F}}
\newcommand{\A}{\mathbb{A}}
\newcommand{\PP}{\mathbb{P}}
\newcommand{\E}{\mathbb{E}}
\newcommand{\eps}{\varepsilon}
\newcommand{\dd}{\,\mathrm{d}}
\begin{document}
\section*{UnsolvedMath ID: 30000529}
\textbf{Source problem code:} catalogue entry
\subsection{Definitions and mathematical statement}
Write $\mathcal C_D(X)$ for the $\infty$-category of $D$-crystals on the scheme $X$ intended by the cited source, and write $K(\mathcal C_D(X))$ for its algebraic $K$-theory spectrum with the source's exact or stable structure. The target, attributed in the supplied assessment to Barwick's Conjecture 17, is
\[
\forall X\text{ a scheme},\qquad
K(\mathcal C_D(X))\simeq K'(X),
\]
with an equivalence natural in $X$ for the class of morphisms and variance intended in the source. An equivalence of spectra induces isomorphisms on all homotopy groups in the degrees of the chosen theory; the question asks for the spectrum-level natural equivalence, not merely an equality of Grothendieck groups.

For a noetherian scheme, the usual interpretation of $K'(X)$ is coherent $K$-theory, also called $G$-theory, built from coherent sheaves. The extract's quantifier is \emph{every scheme}; it supplies neither a general-scheme definition of $K'(X)$ nor the definition, finiteness conditions and exact structure of $D$-crystals. Thus the noetherian interpretation explains common notation but does not replace the stated universal scope by a verified restriction. The symbol $D$ is source-specific and is not identified here with differential operators, ordinary crystals, or arbitrary quasi-coherent complexes.

The supplied literature assessment says that details were not checked beyond the Cohen--Macaulay case. This qualification is retained without asserting that the general equivalence, or even a fully specified theorem for that special case, has been established by the supplied extract.
\subsection{Short English statement}
Ask whether the algebraic $K$-theory spectrum of $D$-crystals is naturally equivalent to coherent $K$-theory for every scheme. The source-specific category, general-scheme conventions and naturality maps are missing, and the reported Cohen--Macaulay limitation remains explicit.
\subsection{Sources}
\begin{itemize}
\item Supplied report: \emph{Algebraic K-Theory} (2007), Oberwolfach Reports, DOI 10.4171/owr/2006/32, pp. 1941--1990; proposed year 2006. \url{https://doi.org/10.4171/owr/2006/32}.
\item Supplied reference: Barwick, \emph{D-crystals}, \url{https://www.maths.ed.ac.uk/~cbarwick/papers/D-crys.pdf}; the registry's assessment identifies Conjecture 17. The linked PDF has not been accessed in preparing this draft.
\item Local UnsolvedMath registry, permanent ID 30000529, catalogue entry; created 25 August 2026, updated 16 September 2026. The corrected\_verified note says the equivalence question is explicit. Its 21 August 2026 literature triage labels the general assertion open and reports that details were not checked beyond the Cohen--Macaulay case. This dated assessment is supplied provenance requiring independent verification.
\end{itemize}
\end{document}
